\documentclass[11pt]{article}
\usepackage[a4paper,margin=27mm]{geometry}
\usepackage{amsmath,amssymb,amsthm,booktabs}
\newtheorem{theorem}{Theorem}
\newtheorem{corollary}[theorem]{Corollary}
\newcommand{\Aut}{\operatorname{Aut}}
\newcommand{\C}{\operatorname{C}}
\title{The order-five row for O'Nan's group}
\author{}\date{}
\begin{document}\maketitle
Let $G$ be O'Nan's simple group and $A=\Aut(G)$. We use the conventions
of Connor and Leemans, \emph{Algorithmic enumeration of regular maps},
Ars Math. Contemp. 10 (2016), 211--222. Thus a pair $(r,s)$ generates
$G$, satisfies $|rs|=2$, and is counted up to simultaneous conjugacy
under $A$. Inner and outer refer to the unique automorphism inverting
both rotations. There is no further division by duality, even when the
rotation orders are equal.

\begin{theorem}
The complete counts for $|r|=5$ and $q=|s|\ge5$ are as follows.
\[
\begin{array}{r|rrrr}
q&\mathrm{all}&\mathrm{inner}&\mathrm{outer}&\mathrm{chiral}\\\hline
5&26&2&2&22\\
6&98&4&4&90\\
7&150&1&15&134\\
8&470&18&0&452\\
10&365&5&14&346\\
11&718&11&17&690\\
12&211&13&6&192\\
14&290&13&13&264\\
15&340&21&15&304\\
16&1966&46&0&1920\\
19&1242&30&44&1168\\
20&800&0&26&774\\
28&560&24&26&510\\
31&502&0&28&474\\\hline
\mathrm{sum}&7738&188&210&7340
\end{array}
\]
\end{theorem}

\begin{proof}
Use the established correspondence with triples $(a,b,c)$ of
involutions in the same coset of $G$, satisfying $ac=ca$ and $a\ne c$.
The rotations are $ab$ and $bc$. Fix the first involution $a$, and put
$J=\C_G(a)$ and $K=\C_A(a)$. The acquired centralizers satisfy
$|K|=2|J|$. The complete order-five partner certificate gives two
$J$-orbits on the inner partners $b$, of lengths $40320$ and $4480$,
and two on the outer partners, both of length $29260$.
Their stabilizers $H_i=\C_J(a,b_i)$ have orders $4,36,6,6$.
These are actual embedded stabilizers with checked word generators;
the orbit cover is not inferred from their abstract names.

The entire same-parity commuting domain for $c$ has size $1750$ in the
inner case and $1463$ in the outer case. Its partitions under these
four stabilizers have respectively $462,74,261,261$ orbits. For every
orbit the literal reader lifts a representative, checks the involution
and commutation identities, computes $|b_i c|$, and decides whether
$\langle ab_i,b_i c\rangle=G$. For this last decision it computes a point
orbit in the acquired permutation representation of degree $245520$.
The socle orbit has size $122760$, whereas none of the thirteen orders
in the exhaustive maximal-subgroup roster is divisible by $122760$.
Hence this orbit size is equivalent to generation: a proper generated
subgroup lies in a maximal subgroup whose order is divisible by the
size of every one of its point orbits. The ambient model and maximal
roster are explicit borrowed premises, as in the accepted order-four
calculation.

Let $M_i(q)$ be the number of eligible third involutions in family $i$
that generate and have second order $q$. Each generating $H_i$-orbit
has size $|H_i|$; this is checked separately for every such orbit.
Indeed, an element stabilizing the triple fixes both generators of $G$
and is therefore the identity automorphism. The same argument makes
the action of $K$ on all generating triples with first entry $a$ free.
Consequently the number of $A$-classes for each parity is
\[
 \frac{1}{|K|}\sum_i\frac{|J|}{|H_i|}M_i(q)
 =\frac12\sum_i\frac{M_i(q)}{|H_i|},
\]
where the sum includes both partner orbits of that parity. This is the
fixed-involution counting theorem already proved with the partner
certificate. In particular no assumption about outer fusion of the two
$J$-orbits is needed, and an individual half-summand need not be integral.

Applying the formula to the four complete journals gives the inner and
outer columns. Subtract them from the accepted total count in each cell
to obtain the chiral column. The calculations also include second order
four: they give zero inner and two outer classes, agreeing with the
separately accepted order-four row. Those entries are excluded only from
the displayed minimum-order-five sums. No expected source total or
expected split is an input to the calculations.
\end{proof}

\begin{corollary}
Type $\{5,5\}$ has two inner-reflexible, two outer-reflexible and
twenty-two chiral classes.
\end{corollary}
This recovers the corresponding split in Section 5 and the order-five
totals in Section 7 of the source. The historical running time is
attributed rather than reproduced. The inversion correspondence,
group and character inputs, accepted numerical table, and complete
commuting domains retain their earlier attribution. The contribution
here is the explicit assembly through the four acquired partner orbits
and its literal certificate; no priority claim is made for the orbit
counting principle. The full inner and outer type spectra, the remaining
order-six chirality claim, and whole-source completion are separate
obligations. No new numerical cell, open conjecture or Lean result is
claimed by this row calculation.
\end{document}
